% GENERATED FILE. Edit canonical lessons, then run npm run generate:books.
% canonical-source-hash: fnv1a64:a638bb9af52b82a3
% canonical-lessons: UR-C335-khilaf, UR-C335-himayat-karna, UR-C335-mukhalafat-karna, UR-C335-muttafiq, UR-C335-fariq

\chapter{Taking Sides}
\label{ch:ur-taking-sides}
\hlchaptermodality{335}

\begin{chapteropening}
\textbf{By the end of this chapter:} \emph{I can say who is for an idea and who is against it.}

Complete the last lesson of chapter 335: I can say who is for an idea and who is against it.
\end{chapteropening}

\section[\texorpdfstring{\ur{خلاف} (khilāf)}{خلاف (khilāf)}]{\ur{خلاف} (khilāf) — against (\ur{کے} \ur{خلاف})}

\label{lesson:UR-C335-khilaf}



\begin{quote}
Before the new one: say the Urdu for an altercation, then the Urdu for logic.
\end{quote}

\subsection*{You'll want to know: \ur{خلاف}}
\textbf{\ur{خلاف}} — \emph{khilāf} — \textquotedblleft{}against (\ur{کے} \ur{خلاف})\textquotedblright{}.

\begin{quote}
\ur{شہری} \ur{اس} \ur{نئے} \ur{قانون} \ur{کے} \ur{خلاف} \ur{ہیں۔}
\end{quote}

\begin{quote}
\emph{shahrī is naye qānūn ke khilāf haiṅ.}
\end{quote}

\begin{quote}
The citizens are against this new law.
\end{quote}

\begin{cousinweb}
\textbf{\ur{خلاف}} (\emph{khilāf}) is usually traced, through Persian, to Arabic \emph{khilāf}, \textquotedblleft{}opposition, the contrary\textquotedblright{}.
\end{cousinweb}

\begin{grammarlens}[title={What you've built}]
One more word for giving reasons and opinions.
\end{grammarlens}

\subsection*{Your turn}
\begin{itemize}
\raggedright
  \item \emph{Say it:} \emph{khilāf}
  \item \emph{Say it:} \emph{khilāf}, once more
  \item \emph{Say it:} \emph{shahrī is naye qānūn ke khilāf haiṅ.} — The citizens are against this new law.
  \item \emph{Recall:} say the Urdu for an altercation, then the Urdu for logic, then say \emph{khilāf} again
\end{itemize}

\begin{tcolorbox}[breakable,colback=teal!4,colframe=teal!35!black,title={Before you move on}]
What does \ur{خلاف} mean? (\textquotedblleft{}Against (\ur{کے} \ur{خلاف})\textquotedblright{}.) Say it once more.
\end{tcolorbox}
\section[\texorpdfstring{\ur{حمایت} \ur{کرنا} (himāyat karnā)}{حمایت کرنا (himāyat karnā)}]{\ur{حمایت} \ur{کرنا} (himāyat karnā) — to back; to support (an idea, a person)}

\label{lesson:UR-C335-himayat-karna}



\begin{quote}
Before the new one: say the Urdu for logic, then the Urdu for against.
\end{quote}

\subsection*{You'll want to know: \ur{حمایت} \ur{کرنا}}
\textbf{\ur{حمایت} \ur{کرنا}} — \emph{himāyat karnā} — \textquotedblleft{}to back; to support (an idea, a person)\textquotedblright{}.

\begin{quote}
\ur{میں} \ur{آپ} \ur{کی} \ur{رائے} \ur{کی} \ur{حمایت} \ur{کرتا} \ur{ہوں۔}
\end{quote}

\begin{quote}
\emph{maiṅ āp kī rāe kī himāyat kartā hūṅ.}
\end{quote}

\begin{quote}
I support your view.
\end{quote}

Past, with \ur{نے} (he or she …): \ur{اس} \ur{نے} \ur{سارا} \ur{کی} \ur{حمایت} \ur{کی}, \emph{us ne sārā kī himāyat kī}.

\begin{cousinweb}
\textbf{\ur{حمایت} \ur{کرنا}} (\emph{himāyat karnā}) is built from \emph{himāyat}, usually traced, through Persian, to Arabic \emph{ḥimāya}, \textquotedblleft{}protection\textquotedblright{}, and \textbf{\ur{کرنا}}, \textquotedblleft{}to do\textquotedblright{}.
\end{cousinweb}

\begin{grammarlens}[title={What you've built}]
One more word for giving reasons and opinions.
\end{grammarlens}

\subsection*{Your turn}
\begin{itemize}
\raggedright
  \item \emph{Say it:} \emph{himāyat karnā}
  \item \emph{Say it:} \emph{himāyat karnā}, once more
  \item \emph{Say it:} \emph{maiṅ āp kī rāe kī himāyat kartā hūṅ.} — I support your view.
  \item \emph{Recall:} say the Urdu for logic, then the Urdu for against, then say \emph{himāyat karnā} again
\end{itemize}

\begin{tcolorbox}[breakable,colback=teal!4,colframe=teal!35!black,title={Before you move on}]
What does \ur{حمایت} \ur{کرنا} mean? (\textquotedblleft{}To back; to support (an idea, a person)\textquotedblright{}.) Say it once more.
\end{tcolorbox}
\section[\texorpdfstring{\ur{مخالفت} \ur{کرنا} (mukhālafat karnā)}{مخالفت کرنا (mukhālafat karnā)}]{\ur{مخالفت} \ur{کرنا} (mukhālafat karnā) — to oppose}

\label{lesson:UR-C335-mukhalafat-karna}



\begin{quote}
Before the new one: say the Urdu for against, then the Urdu for to back.
\end{quote}

\subsection*{You'll want to know: \ur{مخالفت} \ur{کرنا}}
\textbf{\ur{مخالفت} \ur{کرنا}} — \emph{mukhālafat karnā} — \textquotedblleft{}to oppose\textquotedblright{}.

\begin{quote}
\ur{کسان} \ur{نئے} \ur{ٹیکس} \ur{کی} \ur{مخالفت} \ur{کرتے} \ur{ہیں۔}
\end{quote}

\begin{quote}
\emph{kisān naye ṭaiks kī mukhālafat karte haiṅ.}
\end{quote}

\begin{quote}
The farmers oppose the new tax.
\end{quote}

Past, with \ur{نے} (he or she …): \ur{اس} \ur{نے} \ur{قانون} \ur{کی} \ur{مخالفت} \ur{کی}, \emph{us ne qānūn kī mukhālafat kī}.

\begin{cousinweb}
\textbf{\ur{مخالفت} \ur{کرنا}} (\emph{mukhālafat karnā}) is built from \emph{mukhālafat}, usually traced, through Persian, to Arabic \emph{mukhālafa}, \textquotedblleft{}opposing\textquotedblright{}, and \textbf{\ur{کرنا}}, \textquotedblleft{}to do\textquotedblright{}.
\end{cousinweb}

\begin{grammarlens}[title={What you've built}]
One more word for giving reasons and opinions.
\end{grammarlens}

\subsection*{Your turn}
\begin{itemize}
\raggedright
  \item \emph{Say it:} \emph{mukhālafat karnā}
  \item \emph{Say it:} \emph{mukhālafat karnā}, once more
  \item \emph{Say it:} \emph{kisān naye ṭaiks kī mukhālafat karte haiṅ.} — The farmers oppose the new tax.
  \item \emph{Recall:} say the Urdu for against, then the Urdu for to back, then say \emph{mukhālafat karnā} again
\end{itemize}

\begin{tcolorbox}[breakable,colback=teal!4,colframe=teal!35!black,title={Before you move on}]
What does \ur{مخالفت} \ur{کرنا} mean? (\textquotedblleft{}To oppose\textquotedblright{}.) Say it once more.
\end{tcolorbox}
\section[\texorpdfstring{\ur{متفق} (muttafiq)}{متفق (muttafiq)}]{\ur{متفق} (muttafiq) — in agreement}

\label{lesson:UR-C335-muttafiq}



\begin{quote}
Before the new one: say the Urdu for to back, then the Urdu for to oppose.
\end{quote}

\subsection*{You'll want to know: \ur{متفق}}
\textbf{\ur{متفق}} — \emph{muttafiq} — \textquotedblleft{}in agreement\textquotedblright{}.

\begin{quote}
\ur{میں} \ur{آپ} \ur{سے} \ur{متفق} \ur{ہوں۔}
\end{quote}

\begin{quote}
\emph{maiṅ āp se muttafiq hūṅ.}
\end{quote}

\begin{quote}
I agree with you.
\end{quote}

\begin{cousinweb}
\textbf{\ur{متفق}} (\emph{muttafiq}) is usually traced, through Persian, to Arabic \emph{muttafiq}, \textquotedblleft{}agreeing, of one mind\textquotedblright{}.
\end{cousinweb}

\begin{grammarlens}[title={What you've built}]
One more word for giving reasons and opinions.
\end{grammarlens}

\subsection*{Your turn}
\begin{itemize}
\raggedright
  \item \emph{Say it:} \emph{muttafiq}
  \item \emph{Say it:} \emph{muttafiq}, once more
  \item \emph{Say it:} \emph{maiṅ āp se muttafiq hūṅ.} — I agree with you.
  \item \emph{Recall:} say the Urdu for to back, then the Urdu for to oppose, then say \emph{muttafiq} again
\end{itemize}

\begin{tcolorbox}[breakable,colback=teal!4,colframe=teal!35!black,title={Before you move on}]
What does \ur{متفق} mean? (\textquotedblleft{}In agreement\textquotedblright{}.) Say it once more.
\end{tcolorbox}
\section[\texorpdfstring{\ur{فریق} (farīq)}{فریق (farīq)}]{\ur{فریق} (farīq) — a side in a dispute}

\label{lesson:UR-C335-fariq}



\begin{quote}
Before the new one: say the Urdu for to oppose, then the Urdu for in agreement.
\end{quote}

\subsection*{You'll want to know: \ur{فریق}}
\textbf{\ur{فریق}} — \emph{farīq} — \textquotedblleft{}a side in a dispute\textquotedblright{}.

\begin{quote}
\ur{دونوں} \ur{فریق} \ur{ایک} \ur{دوسرے} \ur{سے} \ur{ناراض} \ur{ہیں۔}
\end{quote}

\begin{quote}
\emph{donoṅ farīq ek dūsre se nārāz haiṅ.}
\end{quote}

\begin{quote}
Both sides are angry with each other.
\end{quote}

Masculine.

\begin{cousinweb}
\textbf{\ur{فریق}} (\emph{farīq}) is usually traced, through Persian, to Arabic \emph{farīq}, \textquotedblleft{}a party, a group\textquotedblright{}.
\end{cousinweb}

\begin{grammarlens}[title={What you've built}]
One more word for giving reasons and opinions.
\end{grammarlens}

\subsection*{Your turn}
\begin{itemize}
\raggedright
  \item \emph{Say it:} \emph{farīq}
  \item \emph{Say it:} \emph{farīq}, once more
  \item \emph{Say it:} \emph{donoṅ farīq ek dūsre se nārāz haiṅ.} — Both sides are angry with each other.
  \item \emph{Recall:} say the Urdu for to oppose, then the Urdu for in agreement, then say \emph{farīq} again
\end{itemize}

\begin{tcolorbox}[breakable,colback=teal!4,colframe=teal!35!black,title={Before you move on}]
What does \ur{فریق} mean? (\textquotedblleft{}A side in a dispute\textquotedblright{}.) Say it once more.
\end{tcolorbox}
